\cnabstract

随着制造业数字化转型持续推进，企业内部订单处理、仓储管理、生产执行、采购协同与质量追溯等业务环节之间的联动关系日益紧密。传统以人工沟通和分散系统支撑为主的管理方式，普遍存在信息传递滞后、库存状态不透明、跨部门协同效率较低以及质量问题追踪链条不完整等不足。为提升供应链业务协同能力，本文围绕企业供应链场景，设计并实现了一套基于 Spring Boot 与 Vue 3 的企业供应链协同管理系统。系统采用前后端分离架构，以 Spring Security 与 JWT 完成身份认证和基于角色的权限控制，以 MySQL 与 JPA 实现业务数据持久化，并结合 WebSocket/STOMP 构建实时消息通知机制，从而支持顾客、销售管理员、仓库管理员、生产管理员、采购管理员、供应商、质检管理员以及系统管理员等多角色在统一平台上协同作业。

本文从系统需求分析、总体架构设计、业务流程设计、数据库设计以及关键功能实现等方面展开论述。系统围绕订单驱动主线，构建了销售审核、仓储发货、生产执行、采购供应、质检留痕与质量追溯的闭环业务流程，并在此基础上实现库存预警与周计划生成机制，以增强企业对库存风险和计划衔接问题的响应能力。论文重点分析了角色权限控制、订单至质检协同流程、采购至收货协同流程、实时消息推送、质量追溯以及周计划生成等关键实现机制。测试结果表明，所设计系统能够较好地满足企业供应链协同管理的基本需求，在流程贯通性、数据一致性和业务可追溯性方面具有较好的工程应用价值。

\cnkeyword{供应链协同管理；Spring Boot；Vue 3；库存预警；质量追溯}